% Figure from MATH 452 notes, section 13. Compile with the notes preamble.
\begin{tikzpicture}[scale=1.2]
% ================= left: the (x,y)-plane =================
\begin{scope}
\draw[-stealth, thin] (-1.55,0) -- (1.65,0) node[right, font=\scriptsize] {$x$};
\draw[-stealth, thin] (0,-1.0) -- (0,1.05) node[above, font=\scriptsize] {$y$};
% square Q around (x0,y0) = (1, 0.5)
\fill[figB!12] (0.7,0.25) rectangle (1.3,0.75);
\foreach \a in {0.85,1.0,1.15} \draw[figB!45, thin] (\a,0.25) -- (\a,0.75);
\foreach \b in {0.375,0.5,0.625} \draw[figB!45, thin] (0.7,\b) -- (1.3,\b);
\draw[figB, thick] (0.7,0.25) rectangle (1.3,0.75);
\filldraw (1,0.5) circle (1.1pt);
\node[font=\scriptsize, anchor=south, inner sep=2pt] at (1,0.75) {$(x_0,y_0)=(1,\tfrac12)$};
% antipodal square -Q
\fill[figR!10] (-1.3,-0.75) rectangle (-0.7,-0.25);
\draw[figR, thick, dashed] (-1.3,-0.75) rectangle (-0.7,-0.25);
\filldraw[figR] (-1,-0.5) circle (1.1pt);
\node[font=\scriptsize, figR, anchor=north, inner sep=2pt] at (-1,-0.75) {$(-x_0,-y_0)$};
% the origin, where det J = 0
\filldraw[fill=white, draw=black] (0,0) circle (1.1pt);
\end{scope}
% ================= arrows =================
\draw[-stealth, thick] (1.95,0.75) .. controls (2.35,0.95) .. (2.75,0.75);
\node[font=\scriptsize, anchor=south] at (2.35,0.93) {$(\varphi,\psi)$};
\draw[-stealth, thick, figG] (2.75,0.3) .. controls (2.35,0.1) .. (1.95,0.3);
\node[font=\scriptsize, figG!80!black, anchor=north, align=center] at (2.35,0.12) {$(f,g)$\\ local};
% ================= right: the (u,v)-plane =================
\begin{scope}[shift={(3.45,-0.25)}]
\draw[-stealth, thin] (-0.4,0) -- (1.95,0) node[right, font=\scriptsize] {$u$};
\draw[-stealth, thin] (0,-0.2) -- (0,2.2) node[above, font=\scriptsize] {$v$};
% image region, boundary traced counterclockwise
\fill[figB!12]
  plot[domain=0.7:1.3, samples=20] (\x*\x-0.0625, 0.5*\x)
  -- plot[domain=0.25:0.75, samples=20] (1.69-\x*\x, 2.6*\x)
  -- plot[domain=1.3:0.7, samples=20] (\x*\x-0.5625, 1.5*\x)
  -- plot[domain=0.75:0.25, samples=20] (0.49-\x*\x, 1.4*\x) -- cycle;
\foreach \a in {0.85,1.0,1.15}
  \draw[figB!45, thin, domain=0.25:0.75, samples=20] plot ({\a*\a-\x*\x}, {2*\a*\x});
\foreach \b in {0.375,0.5,0.625}
  \draw[figB!45, thin, domain=0.7:1.3, samples=20] plot ({\x*\x-\b*\b}, {2*\x*\b});
\draw[figB, thick]
  plot[domain=0.7:1.3, samples=20] (\x*\x-0.0625, 0.5*\x)
  -- plot[domain=0.25:0.75, samples=20] (1.69-\x*\x, 2.6*\x)
  -- plot[domain=1.3:0.7, samples=20] (\x*\x-0.5625, 1.5*\x)
  -- plot[domain=0.75:0.25, samples=20] (0.49-\x*\x, 1.4*\x) -- cycle;
\filldraw (0.75,1) circle (1.1pt);
\node[font=\scriptsize, anchor=north west, inner sep=1pt, fill=figB!12] at (0.78,0.97) {$(u_0,v_0)$};
\node[font=\scriptsize, figB, align=left, anchor=west] at (1.2,1.95) {image of $Q$\\ \textcolor{figR}{and of $-Q$}};
\end{scope}
\node[font=\scriptsize, figB, anchor=west] at (1.32,0.5) {$Q$};
\node[font=\scriptsize, figR, anchor=east] at (-1.32,-0.5) {$-Q$};
\end{tikzpicture}
